\chapter{M12: stationäre bidirektionale FSI}\label{vv:fsi}
Dokument 93 legt den elastisch geführten Kanaldeckel aus Dokument 71 fest:
Höhe \(H_0=0.1\) m, Länge und Breite 1 m, fester Volumenstrom über ein
entwickeltes Einlassprofil, Druckauslass. Ein Tet4-Deckel mit dem
\(10^9\)-fachen Modul liegt auf einer einachsigen Federschicht; damit gilt
\(k_s=E_s bL/h_s\) exakt. Für den starren Deckel folgt
\[
d\,(H_0+d)^3=\frac{6\mu L^2 Q}{k_s}.
\]
Im SA-Fall liefert der unabhängige Randwertlöser aus Kapitel~\ref{vv:rans}
den festen Volumenstrom bei \(Re_\tau=587.19\). Bei festem Volumenstrom bleibt
\(Re_\tau\) erhalten, also \(\Delta p\propto H^{-3}\); der Test prüft dies mit
dem Randwertlöser an der finalen Höhe bis \(10^{-6}\).

\begin{center}
\begin{tabular}{lrrrr}
\toprule
Fall & Netzpaare & Fehler fein [Prozent] & letzte Änderung [Prozent] & Effekt [Prozent] \\
\midrule
laminar & 4 & 0.048 & 0.144 & 1.747 (1.748) \\
SA-noft2 & 4 & 0.119 & 0.112 & 2.855 (2.859) \\
\bottomrule
\end{tabular}
\end{center}
Grenzen: laminar 1 und 0.2 Prozent, SA 2 und 0.4 Prozent. Der Effekt ist der
gemessene bidirektionale Unterschied zur Einweglösung, in Klammern der
analytische Wert. Jedes Netzpaar erreicht beide Kopplungsgates in fünf
Schritten; Transferaudit und virtuelle Schnittstellenarbeit bleiben unter
\(3\cdot10^{-15}\), maximales \(y^+\) sinkt von 0.498 auf 0.158. Der Nachweis
gilt für den geführten Deckel in voll entwickelter Kanalströmung und ist keine
experimentelle Validierung. Unabhängiges menschliches Fachreview liegt nicht vor.
